\section{Actual original global-section tensor powers and true model norms}
\label{sec:scheme-global-section-tensor-power-norms}

\begin{definition}[The actual original section power]
\label{def:scheme-global-section-tensor-power}
\lean{EnvelopeProperty.SchemeGlobalSectionTensorPowerNorms.powerSectionHom,EnvelopeProperty.SchemeGlobalSectionTensorPowerNorms.powerSection,EnvelopeProperty.SchemeGlobalSectionTensorPowerNorms.sectionHom_powerSection}
\leanok
\uses{lem:scheme-original-global-section-hom,lem:fc-original-tensor-section-homs}
For an actual original section $s\in\Gamma(X,M)$, form its original
unit-to-module morphism. Its zeroth power is the original identity of the
structure module. Its successor is the actual original tensor of the previous
morphism with the same original section morphism, followed by the genuine
comparison with the original explicit sheafified tensor. Evaluate this original
unit morphism at one to obtain the actual entire section of $M^{\otimes n}$.
The original section-unit equivalence returns exactly this recursive morphism.
\end{definition}
\begin{proof}
A morphism out of the original structure module is determined by its entire
top section at one: original presheaf naturality restricts this value to every
open, and original module linearity determines every scalar input. The
original whole section-unit correspondence therefore inverts evaluation at
one on each recursively formed morphism.
\end{proof}

\begin{theorem}[True original section-power norms]
\label{thm:scheme-global-section-tensor-power-norm}
\lean{EnvelopeProperty.SchemeGlobalSectionTensorPowerNorms.modelHomNorm_tensorSectionHom,EnvelopeProperty.SchemeGlobalSectionTensorPowerNorms.modelHomNorm_powerSectionHom,EnvelopeProperty.SchemeGlobalSectionTensorPowerNorms.globalSectionNorm_powerSection}
\leanok
\uses{def:scheme-global-section-tensor-power,thm:scheme-global-section-model-norm-isometry,lem:actual-model-hom-norm-composition,lem:actual-model-tensor-hom-norm,lem:actual-line-fiber-section-model-norm}
On any actual original proper model over a prescribed bounded base valuation,
let $M$ be any invertible original module and $s$ any entire original global
section. At every actual original native point, the true original proper-model
norm of the actual $n$th section power is $\lVert s\rVert^n$. This includes
zero norms and $n=0$.
\end{theorem}
\begin{proof}
The original tensor section is the inverse original structure multiplication
followed by the tensor of the two section morphisms. Genuine model morphism
norms multiply under composition and under the actual tensor of whole
morphisms. The original structure comparison is an isomorphism and has norm
one. Induction, with the identity at degree zero, proves the claimed norm
formula for every original section-unit power. The actual original fibre
section of a unit-to-line morphism has norm equal to that morphism norm.
The preceding exact original section-unit identity gives the stated formula
for the entire original global-section power.
\end{proof}

\begin{theorem}[Original section powers in the actual exponent block]
\label{thm:scheme-global-section-block-power-norm}
\lean{EnvelopeProperty.SchemeGlobalSectionTensorPowerNorms.blockPowerSection,EnvelopeProperty.SchemeGlobalSectionTensorPowerNorms.globalSectionNorm_blockPowerSection}
\leanok
\uses{thm:scheme-global-section-tensor-power-norm,thm:scheme-global-section-model-norm-isometry,lem:bundle-tensor-power-multiplication}
For every entire original $s\in\Gamma(X,M^{\otimes d})$, transport its actual
$r$th section power through the genuine original whole module isomorphism
$(M^{\otimes d})^{\otimes r}\simeq M^{\otimes dr}$. Its true original model
norm is exactly $\lVert s\rVert^r$ at every actual original native point.
\end{theorem}
\begin{proof}
Apply the preceding theorem to the original invertible module
$M^{\otimes d}$. The actual whole original module isomorphism preserves the
true model norm of every entire section by the genuine isometry theorem.
Compose these two original norm identities. There is no norm, section-lifting
or exponent-block comparison premise.
\end{proof}
